\section{Integrierbare Funktionen}\label{MA2-integrierbare-funktionen}

\begin{table}[H]
  \centering
  \begin{tabularx}{\textwidth}{@{}>{\raggedright\arraybackslash}X r@{}}
    \toprule
    \multicolumn{2}{c}{\textbf{�bersicht �ber Lernziele und Voraussetzungen}}\\
    \midrule
    \textbf{Lernziele} & \\[.5\normalbaselineskip]
    \tabitem ein Skalarprodukt im Funktionenraum definieren und zur Rechnung verwenden & \\
    \tabitem die Norm einer Funktion aus dem Skalarprodukt bestimmen und Funktionen normieren & \\
    \tabitem Orthogonalit�t von Funktionen durch Berechnung des Skalarprodukts pr�fen & \\
    \tabitem den Begriff \emph{Orthonormalsystem} f�r Funktionen formulieren und anwenden & \\
    \tabitem die Orthonormalit�t von Funktionenpaaren bzgl eines Skalarprodukts pr�fen & \\[.5\normalbaselineskip]
    \midrule
    \textbf{Voraussetzungen} & \\[.5\normalbaselineskip]
    \tabitem mit Skalarprodukten und Normen in $\mathbb R^n$ umgehen k�nnen & \abschnittsverweis{MA1-vektoren} \\
    \tabitem Integrale berechnen k�nnen & \abschnittsverweis{MA2-integrationsregeln}\\
    \tabitem trigonometrische Funktionen sicher anwenden k�nnen & \abschnittsverweis{MA1-verkettung-von-funktionen} \\
    \bottomrule
  \end{tabularx}
\end{table}

\noindent \ifintegriereaufgaben Die �bungsaufgaben zum Thema finden Sie \hyperref[MA2-integrierbare-funktionen-aufgaben]{hier}.
\vspace{0.3cm}
\hrule
\fi


\subsection{Vom $\mathbb{R}^n$ zum Funktionenraum}

Im $\mathbb{R}^n$ haben wir das \underline{Skalarprodukt}\randverweis{MA1-vektoren}  kennengelernt, es war wie folgt definiert

\[
	\forall \, \vec x, \vec y \in \mathbb R^n: \quad \langle \vec x,\vec y\rangle = \sum_{i=1}^n x_i \cdot  y_i\,.
\]

Analog definieren wir f�r zwei Funktionen \(f, g\), die �ber dem Intervall $[a,b]$ integrierbar sind\footnote{Stetigkeit auf \((a,b)\) und Beschr�nktheit auf \([a,b]\) garantiert beispielsweise die Integrierbarkeit}:

\[
	\langle f,g\rangle = \int\limits_a^b f(x)\cdot g(x)\,\mathrm d x\, .
\]
Und das ist sinnvoll, denn offensichtlich gilt (genauso wie f�r Spaltenvektoren)
\[
	\langle f,f \rangle = \int\limits_a^b f(x)^2 \,\mathrm d x \geq 0 \, .
\]

Tats�chlich gibt es wenig Unterschied in der Art und Weise wie wir Vektoren miteinander verarbeiten und derjenigen wie wir Funktionen miteinander verarbeiten:
\begin{itemize}
	\item Addition klappt,

\[ f,g \; \quad \textnormal{Funktionen } \Rightarrow f+g \; \quad \textnormal{auch wieder eine Funktion}\,,\]

\item Skalierung klappt
\[ \lambda \textnormal{ Skalar }, f \; \quad \textnormal{Funktion } \Rightarrow \lambda \cdot f \; \quad \textnormal{auch wieder eine Funktion}\,,\]

\item und ein Skalarprodukt haben wir auch!
\end{itemize}
Es wundert einem deshalb kaum, dass wir Begriffe und Strukturen, die wir aus der linearen Algebra kennen, nun wieder auftauchen und zwar zur Analyse von Mengen von Funktionen. Funktionen sind tats�chlich Vektoren, Vektoren in einem unendlichdimensionalen Raum.

Einen Vektor \(\vec x \in \mathbb R^n\) haben wir normiert, in dem wir ihn durch seine L�nge (Norm) geteilt haben, also
\[ \hat  x := \dfrac{\vec x}{\| \vec x\|}\,, \quad \textnormal{ wobei } \;  \| \vec x\| = \sqrt{ \langle \vec x, \vec x \rangle } \,. \]
Und entsprechend l�sst sich aus einer Funktion \(f\) abstrakt dessen normiertes Gegenst�ck defninieren:
\[ \hat  f := \dfrac{f}{\| f\|}\,, \quad \textnormal{ wobei } \;  \| f\| = \sqrt{ \langle f, f \rangle } \,. \]
\begin{itemize}
	\item Das Symbol \(\| f\|\) nennen wir die Norm der Funktion \(f\). Den Begriff L�nge nuzt man f�r Funktionen nicht.
	\item Zur Berechnung von \(\| f\|\) m�ssen wir wissen, wie das Skalarprodukt definiert ist - genauer die Integrationsgrenzen \(a,b\) m�ssen festgelegt sein.
\end{itemize}

\begin{exabox}{Normierung einer Funktion}{ }
	Skalarprodukt:
\[
	\langle f,g \rangle = \int\limits_{-1}^{1} f(x)\cdot g(x) \,\mathrm d x
\]

	Funktion:
	\[ p_1: [-1,1]\to \mathbb R, x \mapsto x \]
		Norm:
		\[ \| p_1 \|^2 = \int\limits_{-1}^1 p_1(x)^2 \, \mathrm d x = \int\limits_{-1}^1 x^2 \, \mathrm d x = \Big[ \frac13 x^3\Big]_{-1}^1= \frac23 \Rightarrow \|p_1\| = \sqrt{\frac23} \,.\]
		Normierung:
		\[ \hat p_1 (x) = \frac{1}{ \|p_1\|} \cdot p_1(x) = \frac{1}{\sqrt{\frac23}} \cdot x = \frac32 \cdot x\,.\]
\end{exabox}


\subsection{Orthogonalit�t von Funktionen}

Das Skalarprodukt ist so wertvoll, weil es Orthognoalit�t definiert: zwei Funktionen hei�en \textbf{orthogonal}, wenn

\[
\langle f,g\rangle = 0 \; \Rightarrow \; f\perp g \,.
\]

Wir schauen uns jetzt Beispiele f�r Funktionenp�rchen und Skalarprodukte an und pr�fen, ob die Funktionen orthogonal sind.

\begin{exabox}{Orthogonale Funktionen}{ }
	Funktionen:
	\[\begin{cases}  f: & [-1,1]\to \mathbb R, x \mapsto 1  \\  g: & [-1,1]\to \mathbb R, x \mapsto x \end{cases}\]
		Skalarprodukt:
		\[ \langle f, g \rangle  = \int\limits_{-1}^1 f(x) \cdot g(x) \, \mathrm d x = \int\limits_{-1}^1 1\cdot x \, \mathrm d x = \Big[ \frac12 x^2\Big]_{-1}^1= \frac12\cdot (1^2- (-1)^2) = 0\]
		Schlussfolgerung: \(f \perp g\)
\end{exabox}


\begin{exabox}{Orthogonale Funktionen}{ }
	Funktionen:
	\[\begin{cases}  f: & [-1,1]\to \mathbb R, x \mapsto 1  \\  g: & [-1,1]\to \mathbb R, x \mapsto x^2 \end{cases}\]
		Skalarprodukt:
	\[ \langle f,  g \rangle  = \int\limits_{-1}^1 f(x) \cdot g(x) \, \mathrm d x = \int\limits_{-1}^1 1\cdot x^2 \, \mathrm d x = \Big[ \frac13 x^3\Big]_{-1}^1= \frac13\cdot (1)^3 - (-1)^3) = \frac23\]
		Schlussfolgerung: \(f\not \perp g\)
\end{exabox}

Monome sind also im Allgemeinen nicht paarweise orthogonal! Und noch ein Beispiel mit harmonischen Schwingungen an - beachte, dass das Integrationsgebiet in diesem Skalarprodukt das Intervall \( [-\pi, \pi]\) ist.

\begin{exabox}{Orthogonale Funktionen}{ }
	Funktionen \(f, g\):
	\[\begin{cases}  f: & [-\pi,\pi]\to \mathbb R, x \mapsto \sin(x)  \\  g: & [-\pi,\pi]\to \mathbb R, x \mapsto \cos(x) \end{cases}\]
		Skalarprodukt:
		\[ \langle f,  g \rangle  = \int\limits_{-\pi}^\pi \sin(x) \cdot \cos(x) \, \mathrm d x = \Big[ \frac12 \sin(x)^2\Big]_{-\pi}^\pi= \frac12\cdot (0- 0) = 0\]
		Schlussfolgerung: \(f \perp g\)
\end{exabox}

Interessanterweise kann man zeigen, dass alle harmonischen Schwingungen der Form \[ s_n: x\mapsto \sin(n\cdot x), \quad n \in \mathbb N^+ \] und \[ c_m: x\mapsto \cos(m\cdot x), \quad m \in \mathbb N_0 \]  �ber \([-\pi,\pi]\) \textbf{paarweise orthogonal} sind. Auf mathematisch gilt also folgendes

\[
	\langle s_n, c_m \rangle = \int\limits_{-\pi}^\pi \sin(n\cdot x) \cdot \cos(m\cdot x)\,\mathrm d x  = 0\,.
\]

Und das ist nicht genug, es gilt auch:

\[
	\begin{array}{rcl}
		\langle s_n, s_m \rangle &=& \displaystyle \int\limits_{-\pi}^\pi \sin(n\cdot x) \cdot \sin(m\cdot x)\,\mathrm d x  = 0, \quad n\not=m\,, \\[2ex]
	\langle c_n, c_m \rangle &=& \displaystyle \int\limits_{-\pi}^\pi \cos(n\cdot x) \cdot \cos(m\cdot x)\,\mathrm d x  = 0, \quad n\not=m\,. \end{array}
\]

Die Funktionen
\[
	\left\{ 1, \sin(n\cdot x), \cos(n\cdot x), \quad n \in \mathbb N\right\}
\]
sind bez�glich des Skalarprodukts
\[
\langle f,g\rangle
=
\int\limits_{-\pi}^{\pi} f(x)\cdot g(x)\,\mathrm dx
\]
orthogonal.

\subsection{Orthonormalsystem}

Eine Menge von Funktionen hei�t \emph{Orthogonalsystem}, wenn die Funktionen normiert und paarweise orthogonal zueinander sind.
Mit dem Skalarprodukt
\begin{equation}\label{skp_orthonormalsystem}
\langle u,v\rangle
=
\frac{1}{\pi}
\int_{1}^{2\pi} u(t)\cdot v(t)\,\mathrm dt
\end{equation}
werden die harmonischen Schwingungen aus dem vorigen Abschnitt automatisch normiert. Wir kennen also ein konkretes Beispiel f�r ein Orthogonalsystem.

Um orthographisch kenntlich zu machen, dass wir die harmonischen Schwingungen ab jetzt immer mit dem Skalarprodukt betrachen, in dem sie ein Orthonormalsystem bilden, nutzen den kleinen Hut:

\[ \left\{ \begin{array}{rcl l} \hat s_n(t): &=& \sin(n\cdot t)\,, &\quad n\geq 1\,, \\  \hat c_n(t): &=& \cos(n\cdot t)\,, & \quad n\geq 1 \,.\end{array}\right.  \]

\begin{exabox}{Orthonormalsystem}{ }
	Die Funktionen
\[
	M = \left\{ 1, \hat s_n, \hat c_n, \; n \in \mathbb N\right\}
\]
bilden bez�glich \eqref{skp_orthonormalsystem} ein Orthonormalsystem und das hei�t, dass f�r alle \(n,m\in \mathbb N^+\) gilt
\[
	\begin{array}{rl}
		\textnormal{paarweise orthogonal:} & \quad \left\{
	\begin{array}{rcl l}
		\langle 1, \hat c_n\rangle &=& 0 & \\
		\langle 1, \hat s_n\rangle &=& 0 & \\
		\langle \hat c_n, \hat s_m\rangle &=& 0\,, &  \\
		\langle \hat c_n, \hat c_m\rangle &=& 0\,, & n\not=m \\
		\langle \hat s_n, \hat s_m\rangle &=& 0\,, & n\not=m \\
\end{array} \right. \\
		\textnormal{normiert} & \quad \left\{
	\begin{array}{rcl}
		\langle 1, 1\rangle &=& 1\\
		\langle \hat s_n, \hat s_n\rangle &=& 1\\
		\langle \hat c_n, \hat c_n\rangle &=& 1
\end{array} \right.
\end{array}
\]

\end{exabox}

Das Orthonormalsystem $M$ besitzt zwei zentrale Eigenschaften:

\begin{itemize}
  \item Alle Funktionen sind $2\pi$-periodisch.
  \item Linearkombinationen bleiben ebenfalls $2\pi$-periodisch.
\end{itemize}

Die gemeinsame Periode ist entscheidend:
�berlagert man zwei Funktionen aus $M$,
so entsteht wieder eine $2\pi$-periodische Funktion.
Dasselbe gilt f�r jedes skalare Vielfache.

Damit verh�lt sich $M$ strukturell wie eine Basis im
endlichdimensionalen Vektorraum �
nur dass wir uns hier in einem Funktionenraum befinden.

In der Fourier-Analyse werden wir dieses Orthonormalsystem
verwenden, um $2\pi$-periodische Funktionen als
Reihenentwicklungen in harmonischen Schwingungen darzustellen.

